% TOP_PROBLEM: {"problemId": "problem.quillen-s-conjecture-on-posets-of-elementary-abelian-p-subgroups", "canonicalTitle": "Quillen's Conjecture on Posets of Elementary Abelian p-Subgroups", "releaseRank": 497, "primaryDomain": "algebra_representation_category", "primaryDomainLabel": "Algebra, representation & category theory", "displayStatus": "open_with_solved_subcases", "releaseStatus": "open_with_solved_subcases"}
% !TeX program = lualatex
% ATTEMPT_STATUS: unresolved
% Model: GPT-6 Astra Ultra; continuation dated 27 September 2026.
\documentclass[11pt]{article}
\usepackage[margin=25mm]{geometry}
\usepackage{fontspec}
\setmainfont{Latin Modern Roman}
\usepackage{amsmath,amssymb,mathtools}
\usepackage{unicode-math}
\setmathfont{Latin Modern Math}
\usepackage{xurl,hyperref}
\hypersetup{colorlinks=true,urlcolor=blue}
\setcounter{secnumdepth}{0}
\setlength{\parindent}{0pt}
\setlength{\parskip}{5pt}
\setlength{\emergencystretch}{3em}
\begin{document}
\section*{\texorpdfstring{Quillen's Conjecture on Posets of Elementary Abelian p-Subgroups}{Quillen's Conjecture on Posets of Elementary Abelian p-Subgroups}}\label{497}
\subsection{Definitions and mathematical statement}
For a finite group $G$ and prime $p$, let $\mathcal A_p(G)$ be the partially ordered set, under inclusion, of nontrivial elementary abelian $p$-subgroups, that is, subgroups isomorphic to $({\mathbb{Z}}/p{\mathbb{Z}})^r$ for some $r\ge1$. Its order complex has these subgroups as vertices and finite strict chains as simplices.

Let $O_p(G)$ be the largest normal $p$-subgroup. Quillen's assertion is
\[
 |\mathcal A_p(G)|\text{ contractible}
                   \quad\Longrightarrow\quad O_p(G)\ne1.
\]
Contractible means homotopy equivalent to a point; it is stronger than having vanishing reduced homology. The empty complex, which occurs when $p\nmid|G|$, is not contractible under this convention and therefore supplies no counterexample.
\par\smallskip\textit{Formulation and concept references:} \href{https://eprints.lancs.ac.uk/id/eprint/158837/1/dm_rev9.pdf}{[S1]}.

\subsection{Short English statement}
If the complex of nontrivial elementary abelian p-subgroups contracts to a point, must the group contain a nontrivial normal p-subgroup?
\subsection{Research attempt}

\paragraph{Scope, source audit, and choice of attack (27 September 2026).}
The prime $p$ and the finite group $G$ are arbitrary; the implication concerns
ordinary contractibility of the entire order complex. A subgroup normalized by
a Sylow subgroup is not necessarily normal in $G$. Likewise a nonempty fixed
point set for each Sylow subgroup does not provide a common fixed point.
These distinctions are central to the attempt below.

The supplied manuscript [S1], Introduction, states the conjecture and records
Quillen's established rank-at-most-two case. The proof of that case below is a
self-contained reconstruction, not a claim of a new subcase. More recent
work of Piterman--Smith extends the Aschbacher--Smith theorem to the remaining
odd primes under the theorem's component hypotheses; its abstract does not
assert an unrestricted solution at all odd primes [E1]. Piterman's 2021 paper
already proves a stronger homological version for groups of $p$-rank at most
four [E2]. Thus the elementary rank-two argument is a diagnostic benchmark,
not the current research frontier. Only these specific source statements are
used; no classification theorem is assumed in the calculations below.

The first route is to turn contractibility into a conjugation-fixed vertex.
This succeeds when the complex is a tree. The second route attempts to glue
the fixed vertices supplied by $p$-local information. A third, disproof-oriented
route tests whether Euler characteristic or deletion of leaves can distinguish
the desired complexes from ordinary contractible complexes. Explicit finite
examples expose the limitations of both proposed shortcuts.

\paragraph{Basic homotopy mechanism, with the normal-core direction checked.}
Write $\mathcal S_p(G)$ for the poset of all nontrivial $p$-subgroups, and put
\[
 r_p(G)=\max\{r:(C_p)^r\leq G\}.
\]
In the absence of a nontrivial elementary abelian subgroup, take $r_p(G)=0$.
Strict subgroup inclusion increases rank, so $|\mathcal A_p(G)|$ has dimension
$r_p(G)-1$ when it is nonempty.

We use the following elementary order-homotopy fact. If $f,g:X\to Y$ are
order-preserving maps of finite posets and $f(x)\leq g(x)$ for every $x$, then
their realizations are homotopic. Indeed the map on the product poset
$X\times\{0<1\}$ given by $f$ at zero and $g$ at one is order-preserving:
for $x\leq x'$ the mixed comparison is $f(x)\leq g(x')$.
The usual prism triangulation realizes this map as a homotopy. In particular,
a zigzag of pointwise comparable order maps connecting the identity to a
constant map contracts the order complex.

\textbf{Lemma 497.1 (a finite-poset fiber argument).}
If $f:X\to Y$ is order-preserving and every poset
$X_y=\{x:f(x)\leq y\}$ is nonempty and contractible, then $|f|$ is a
homotopy equivalence.

\emph{Proof.}
Form a poset $M=X\sqcup Y$ with the original orders and the additional
relations $x\leq y$ exactly when $f(x)\leq y$; there are no relations
from a member of $Y$ to a member of $X$. The map $r:M\to Y$ restricting
to $f$ on $X$ and to the identity on $Y$ satisfies
$\operatorname{id}_M\leq i_Yr$. Hence $i_Y$ is a homotopy equivalence.
Remove the vertices of $Y$ from $M$ in a linear extension of its order,
starting with minimal vertices. When $y$ is removed, its strict lower
poset is precisely $X_y$, since every smaller member of $Y$ has already
been removed. Its link is the join of $|X_y|$ with its remaining strict
upper order complex, and is contractible. If there is no upper vertex,
the link is just $|X_y|$.

Deleting a vertex with a nonempty contractible link preserves homotopy
type: the original complex is obtained from the deletion by attaching the
cone on the link along its base. The base inclusion into this cone is a
homotopy equivalence and a simplicial cofibration, so the attachment has
the same homotopy type as the deletion. This follows also by extending a
contraction of the link over its cone. Successive deletions show that
$i_X:X\to M$ is a homotopy equivalence. Since $r i_X=f$, the assertion
follows.\hfill$\square$

For a nontrivial finite $p$-group $Q$, its center is nontrivial by the class
equation. The subgroup
\[
 E_Q=\{z\in Z(Q):z^p=1\}
\]
is a nontrivial elementary abelian group: the center is abelian, so products
and inverses preserve the exponent-$p$ condition. Every
$A\in\mathcal A_p(Q)$ commutes with $E_Q$, and $AE_Q$ is elementary
abelian. Consequently
\[
 A\ \leq\ AE_Q\ \geq\ E_Q
\]
is a contraction of $|\mathcal A_p(Q)|$. Apply Lemma 497.1 to the inclusion
$\mathcal A_p(G)\hookrightarrow\mathcal S_p(G)$: its lower fiber at $Q$
is exactly $\mathcal A_p(Q)$. This proves the homotopy equivalence of the
two subgroup complexes used here without assuming the conjecture.

If $N\triangleleft G$ is a nontrivial $p$-subgroup, then $QN$ is a
$p$-subgroup for every $Q\in\mathcal S_p(G)$; normality is what makes
the product a subgroup, and
$|QN|=|Q||N|/|Q\cap N|$ is a power of $p$. The zigzag
\[
 Q\ \leq\ QN\ \geq\ N
\]
contracts $|\mathcal S_p(G)|$ and therefore also $|\mathcal A_p(G)|$.
This proves the known direction
$O_p(G)\ne1\Rightarrow |\mathcal A_p(G)|$ contractible. Trying to use
$A\mapsto AN$ directly on $\mathcal A_p(G)$ would be incorrect:
$AN$ need not be abelian or have exponent $p$.

\paragraph{A complete benchmark: contractibility in rank at most two.}
\textbf{Proposition 497.2.}
If $r_p(G)\leq2$ and $|\mathcal A_p(G)|$ is contractible, then
$G$ has a nontrivial normal elementary abelian $p$-subgroup.

\emph{Proof.}
The empty complex is excluded by the hypothesis. In rank one the complex
is a nonempty discrete set. Contractibility forces exactly one vertex,
which conjugation fixes. Suppose now the rank is two. The complex is a
finite graph: vertices have ranks one or two, and edges express inclusion
of a rank-one group in a rank-two group. No higher simplex exists.
Contractibility implies connectedness and $H_1=0$. A connected finite
graph has first Betti number $e-v+1$: delete an edge from each cycle
until a spanning tree remains, with each deleted edge accounting for one
independent cycle. Thus the graph is a tree.

A finite tree has a center consisting of one vertex or the endpoints of
one edge. One proof is to remove all leaves simultaneously whenever more
than two vertices remain. For a tree with at least three vertices the
remaining graph is a nonempty tree, since the internal vertices of a path
between surviving vertices also survive. The process ends at a vertex or
an edge, and each round is invariant under every tree automorphism.
Conjugation by $G$ preserves this center. If the center is a vertex, the
corresponding elementary abelian subgroup is normal. If the center is
an edge, its endpoints have different ranks, and conjugation preserves
subgroup orders; hence it cannot exchange the endpoints. Each endpoint
is therefore fixed individually and again gives a nontrivial normal
elementary abelian subgroup.\hfill$\square$

The same proof applies whenever the order complex has a $G$-invariant
subcomplex that is a nonempty tree: no retraction is needed for that
sufficient condition. The obstacle is constructing such a tree from the
assumed contractibility in larger rank. An arbitrary spanning tree of the
one-skeleton will usually not be invariant under conjugation.

\paragraph{Local fixed points exist even when the global complex is bad.}
\textbf{Lemma 497.3 (explicit contraction of a local fixed poset).}
Let $P\leq G$ be any nontrivial $p$-subgroup. The poset
\[
 \mathcal A_p(G)^P=\{A\in\mathcal A_p(G):uAu^{-1}=A
                       \text{ for every }u\in P\}
\]
has contractible order complex.

\emph{Proof.}
Set $E=E_P\leq Z(P)$ as above. If $P$ normalizes $A$, conjugation acts
on the set of elements of $A$. Every nonfixed orbit has size divisible by
$p$, so $|C_A(P)|\equiv |A|\equiv0\pmod p$. The identity is fixed;
hence there are at least $p$ fixed elements and $B(A)=C_A(P)$ is
nontrivial. It is an elementary abelian subgroup, and $A\mapsto B(A)$
preserves inclusion. Since $B(A)$ centralizes $P$, it commutes with $E$.
Thus $B(A)E$ is elementary abelian and is itself centralized by $P$.
All maps in
\[
 A\ \geq\ B(A)\ \leq\ B(A)E\ \geq\ E
\]
are maps of the displayed fixed poset, are order-preserving, and have the
indicated pointwise comparisons. The order-homotopy fact contracts its
realization.\hfill$\square$

There is no use of global contractibility in Lemma 497.3. Moreover the
topological fixed point set of the conjugation action of $P$ on
$|\mathcal A_p(G)|$ equals $|\mathcal A_p(G)^P|$. A point belongs to a
unique open simplex; if fixed, that simplex is stabilized. Its chain has
at most one subgroup of each rank, so a stabilizing conjugation fixes
each vertex. The reverse containment is immediate.

Counting orbits of simplices under this $p$-group action gives the useful
but weak consequence
\begin{equation}\label{ra497:euler}
 \chi(|\mathcal A_p(G)|)\equiv
 \chi(|\mathcal A_p(G)^P|)=1\pmod p.
\end{equation}
Indeed every nonfixed orbit of simplices has cardinality divisible by
$p$, in each dimension separately. The congruence therefore survives the
alternating sum. This shows why a modular Euler-characteristic test does
not recognize contractibility: the same congruence is automatic for every
group with a nontrivial $p$-subgroup. One must retain integral homology,
connectivity, or more structure instead of only this residue.

The failed gluing step can now be written precisely. For the Sylow
subgroups $P_1,\ldots,P_t$, the local fixed complexes are each
contractible, but the desired fixed vertices lie in
\[
 \bigcap_{i=1}^t\mathcal A_p(G)^{P_i}
       \ \cap\ \mathcal A_p(G)^{\langle\text{remaining conjugations}\rangle}.
\]
The expression stresses that Sylow $p$-subgroups need not generate $G$.
Even pairwise overlap of local contractions has not been established,
and a chosen Sylow's central subgroup $E_P$ is generally moved by $G$.
There is no compactness or averaging operation on elementary abelian
subgroups that turns these local vertices into a global one. In particular,
the join of two elementary abelian subgroups need not be a $p$-group.

\paragraph{Adversarial examples and a reproducible finite computation.}
The accompanying standard-Python script enumerates the permutations in
$S_3,A_4,S_4,A_5$ and enumerates elementary abelian subgroups by adjoining
commuting elements of order $p$. This enumeration is complete: any
elementary abelian subgroup has a basis, and adjoining that basis constructs
it along this search. It then counts all strict chains, components of the
one-skeleton, and the subgroups fixed by every conjugation. It never infers
contractibility merely from Euler characteristic. Selected output is
\[
\begin{array}{c|c|c|c|c}
 (G,p)&(f_0,f_1,f_2)&\text{components}&\chi&
       \text{orders of normal elementary subgroups}\\ \hline
 (S_3,2)&(3,0,0)&3&3&\text{none}\\
 (A_4,2)&(4,3,0)&1&1&4\\
 (S_4,2)&(13,12,0)&1&1&4\\
 (A_5,2)&(20,15,0)&5&5&\text{none}\\
 ((C_2)^3,2)&(15,35,21)&1&1&2,4,8
\end{array}
\]
In the last row the final column lists the distinct orders, not the number
of subgroups of each order. The full output also tests $p=3$ in the four
permutation groups and $p=5$ in $A_5$.

The examples can be checked without trusting the program. In $S_3$ at
$p=2$ there are three isolated order-two subgroups. In $A_5$ at $p=2$,
each of the five Klein four groups contains three double transpositions;
distinct such Klein four groups have trivial intersection. Its complex is
therefore five disjoint three-leaf stars. Every local fixed complex from
Lemma 497.3 still contracts, but there is no globally normal elementary
abelian subgroup. This is a counterexample to local-to-global gluing in
the absence of the conjecture's global hypothesis, not to Quillen's
conjecture itself.

The rank-three example is a sharper test of the leaf-removal strategy.
For $V=(C_2)^3$, the nonzero subspaces consist of seven lines, seven
planes, and $V$. There are $21$ line--plane incidences, seven line--$V$
incidences and seven plane--$V$ incidences. The $21$ flags
$\ell<W<V$ are the triangles. Its order complex is a cone because it has
a maximum vertex $V$, so it is certainly contractible. Yet the one-skeleton
has minimum degree four and first Betti number
\[
 35-15+1=21.
\]
There is no leaf to delete. The $21$ two-simplices kill the cycles, a
mechanism absent in Proposition 497.2. Thus replacing ``contractible
complex'' by ``tree one-skeleton'' silently loses valid instances at the
first higher rank. The cone has a fixed vertex for a different reason;
its failure to have leaves is not a counterexample to the target.

\paragraph{Outcome and the first remaining gap.}
\textbf{Outcome: UNRESOLVED.} The work proves the normal-core direction,
reconstructs the established rank-at-most-two converse, and supplies an
explicit contraction for every local $p$-subgroup fixed complex. The
finite checks corroborate these arguments and disprove two intermediate
shortcuts: local fixed points alone do not glue, and the one-skeleton of
a contractible subgroup complex need not be a tree.

The first missing statement for the chosen approach is a mechanism that,
assuming $|\mathcal A_p(G)|$ contractible, produces a nonempty
$G$-invariant tree subcomplex or otherwise a conjugation-fixed elementary
abelian subgroup. No such mechanism has been proved. The exact global
fixed-vertex conclusion is equivalent to the target: a fixed vertex is a
normal $p$-subgroup, while a nontrivial $O_p(G)$ has the characteristic
nontrivial elementary abelian subgroup $E_{O_p(G)}$.

Concrete next tests are to study how rank-three and higher simplices could
be removed in conjugation orbits, to inspect intersections of the local
fixed complexes rather than their individual contractions, and to seek an
integral homology obstruction when $O_p(G)=1$. Existing results cited above
already go beyond the low-rank benchmark, so none of these observations
is advertised as a new frontier theorem. This notebook neither proves nor
disproves the unrestricted assertion.

\subsection{Sources}
\begin{itemize}
\item[S1]A. Díaz Ramos, N. Mazza, A geometric approach to Quillen's conjecture, accepted manuscript, Lancaster University EPrints 158837.
\url{https://eprints.lancs.ac.uk/id/eprint/158837/1/dm_rev9.pdf}.
\textit{Formulation source.}
\item[E1] K. I. Piterman and S. D. Smith,
\emph{Some results on Quillen's conjecture via equivalent-poset techniques},
Journal of Combinatorial Algebra 9 (2025), 265--387, abstract and Introduction.
\url{https://ems.press/content/serial-article-files/51406}.
\textit{Primary published source; consulted 27 September 2026.}
\item[E2] K. I. Piterman,
\emph{An approach to Quillen's conjecture via centralisers of simple groups},
Forum of Mathematics, Sigma 9 (2021), e48, abstract and introductory results.
\url{https://www.cambridge.org/core/journals/forum-of-mathematics-sigma/article/an-approach-to-quillens-conjecture-via-centralisers-of-simple-groups/8DB53651BB6EB02306CA4BB831E9F576}.
\textit{Primary published source; consulted 27 September 2026.}
\end{itemize}
\end{document}
